\documentclass[10pt,a4paper]{amsart}
\usepackage{amsmath,amssymb,amsthm,graphicx}
\usepackage{verbatim}
\newenvironment{reference}{\comment}{\endcomment}
%\newenvironment{reference}{}{}
\newenvironment{slogan}{\comment}{\endcomment}
\newenvironment{history}{\comment}{\endcomment}

% For commutative diagrams we use Xy-pic
\usepackage[all]{xy}

% We use 2cell for 2-commutative diagrams.
\xyoption{2cell}
% Theorem environments.
%
\theoremstyle{plain}
\newtheorem{theorem}[subsection]{Theorem}
\newtheorem{proposition}[subsection]{Proposition}
\newtheorem{lemma}[subsection]{Lemma}

\theoremstyle{definition}
\newtheorem{definition}[subsection]{Definition}
\newtheorem{example}[subsection]{Example}
\newtheorem{exercise}[subsection]{Exercise}
\newtheorem{situation}[subsection]{Situation}

\theoremstyle{remark}
\newtheorem{remark}[subsection]{Remark}
\newtheorem{remarks}[subsection]{Remarks}

\numberwithin{equation}{subsection}

% Macros
%
\def\lim{\mathop{\mathrm{lim}}\nolimits}
\def\colim{\mathop{\mathrm{colim}}\nolimits}
\def\Spec{\mathop{\mathrm{Spec}}}
\def\Hom{\mathop{\mathrm{Hom}}\nolimits}
\def\Ext{\mathop{\mathrm{Ext}}\nolimits}
\def\SheafHom{\mathop{\mathcal{H}\!\mathit{om}}\nolimits}
\def\SheafExt{\mathop{\mathcal{E}\!\mathit{xt}}\nolimits}
\def\Sch{\mathit{Sch}}
\def\Mor{\mathop{\mathrm{Mor}}\nolimits}
\def\Ob{\mathop{\mathrm{Ob}}\nolimits}
\def\Sh{\mathop{\mathit{Sh}}\nolimits}
\def\NL{\mathop{N\!L}\nolimits}
\def\CH{\mathop{\mathrm{CH}}\nolimits}
\def\proetale{{pro\text{-}\acute{e}tale}}
\def\etale{{\acute{e}tale}}
\def\QCoh{\mathit{QCoh}}
\def\Ker{\mathop{\mathrm{Ker}}}
\def\Im{\mathop{\mathrm{Im}}}
\def\Coker{\mathop{\mathrm{Coker}}}
\def\Coim{\mathop{\mathrm{Coim}}}

% Boxtimes
%
\DeclareMathSymbol{\boxtimes}{\mathbin}{AMSa}{"02}

%
% Macros for moduli stacks/spaces
%
\def\QCohstack{\mathcal{QC}\!\mathit{oh}}
\def\Cohstack{\mathcal{C}\!\mathit{oh}}
\def\Spacesstack{\mathcal{S}\!\mathit{paces}}
\def\Quotfunctor{\mathrm{Quot}}
\def\Hilbfunctor{\mathrm{Hilb}}
\def\Curvesstack{\mathcal{C}\!\mathit{urves}}
\def\Polarizedstack{\mathcal{P}\!\mathit{olarized}}
\def\Complexesstack{\mathcal{C}\!\mathit{omplexes}}
% \Pic is the operator that assigns to X its picard group, usage \Pic(X)
% \Picardstack_{X/B} denotes the Picard stack of X over B
% \Picardfunctor_{X/B} denotes the Picard functor of X over B
\def\Pic{\mathop{\mathrm{Pic}}\nolimits}
\def\Picardstack{\mathcal{P}\!\mathit{ic}}
\def\Picardfunctor{\mathrm{Pic}}
\def\Deformationcategory{\mathcal{D}\!\mathit{ef}}
\usepackage{cleveref}
\crefname{originaldefinition}{original definition}{original definitions}
\Crefname{originaldefinition}{Original definition}{Original definitions}
\makeatletter
\newlabel{algebra-definition-formally-smooth}{{00TI}{source}{}{sourcepointer-1}{}}
\newlabel{algebra-definition-formally-smooth@cref}{{[originaldefinition][0][]00TI}{[1][0][]source}{}{}{}}
\makeatother
\crefname{originalreference}{original reference}{original references}
\Crefname{originalreference}{Original reference}{Original references}
\makeatletter
\newlabel{lemma-surjective-cotangent-space}{{06GZ}{source}{}{sourcepointer-2}{}}
\newlabel{lemma-surjective-cotangent-space@cref}{{[originalreference][0][]06GZ}{[1][0][]source}{}{}{}}
\makeatother
\crefname{originalreference}{original reference}{original references}
\Crefname{originalreference}{Original reference}{Original references}
\makeatletter
\newlabel{remark-compare-smooth-schlessinger}{{06HJ}{source}{}{sourcepointer-3}{}}
\newlabel{remark-compare-smooth-schlessinger@cref}{{[originalreference][0][]06HJ}{[1][0][]source}{}{}{}}
\makeatother
\title{A morphism of formal spectra is smooth exactly when the corresponding ring is a power series ring}
\author{The Stacks Project Authors}
\date{stacks-project:formal-defos.tex:local-snapshot-20261012; GNU FDL 1.2 or later}
\begin{document}
\maketitle
\noindent\textit{Editorial note.} All delivered mathematical text is the author's own, copied from the pinned Stacks Project chapter file formal-defos.tex; no statement or proof step was rewritten, reordered or removed.\par\medskip
\noindent\textit{Editorial note.} The counted proof is the author's complete argument: the successive lifting of the map to the quotients by powers of the maximal ideal and the construction of the section compared through relative cotangent spaces, and the explicit construction of the dotted arrow from lifts of the images of the variables in the converse direction.\par\medskip
\noindent\textit{Editorial note.} The label lemma-smooth-morphism-power-series names the counted statement, and the proof environment immediately following it in the source is delivered verbatim as its proof.\par\medskip
\noindent\textit{Editorial note.} The counted statement is carried by the author's own lemma environment in the chapter source.\par\medskip
\noindent\textit{Editorial note.} Pre-open gate: python operations/agent-authoring/tools/precard\_gate.py --source formal-defos.tex --label lemma-smooth-morphism-power-series --start 2241 --end 2319 --local-id 28013 returned RESULT=OK (L28013 free, occupancy FREE, public interval clear). Candidate taken from sources/stacks-formal-defos/candidates-halfA.json.\par\medskip
\noindent\textit{Editorial note.} The statement and proof are copied from the pinned chapter file; the cited inputs are kept as external pointers with their canonical Stacks tags.\par\medskip
\noindent\textit{Allowed standard prerequisites (editorial).}\begin{itemize}
\item Reformulation of smoothness for the functors attached to the rings. Same chapter, the remark comparing smoothness of the associated functors with the lifting property for surjections in the category of Artinian algebras. Supplies the lifting formulation of smoothness used in both directions.
\item Surjectivity of a local map inducing an isomorphism on relative cotangent spaces. Same chapter, the lemma that a map in the completed category which induces an isomorphism on relative cotangent spaces is surjective. Supplies the conclusions that the constructed map f and later the section s are isomorphisms.
\item The lifting property along the quotients by powers of the maximal ideal. Standard consequences of smoothness for the category of Artinian algebras, applied successively to the quotients by powers of the maximal ideal. Supplies the successive lifts f\_2, f\_3, ... whose limit gives the map f : S -> T.
\end{itemize}\medskip
\ifdefined\hypertarget\hypertarget{sourcepointer-1}{}\fi
\noindent\textit{Source reference.} Original definition 00TI: sources/stacks-chow/tags.txt, source lines 1076--1076. Prerequisite pointer invoked at proof line(s) 2265; it names a standard fact of a sibling chapter that the counted argument uses, and every step of the counted proof built on it is retained verbatim in the delivered body.\par\medskip
\ifdefined\hypertarget\hypertarget{sourcepointer-2}{}\fi
\noindent\textit{Source reference.} Original reference 06GZ: sources/stacks-formal-defos/pinned-original-source.tex, source lines 1086--1086. Prerequisite pointer invoked at proof line(s) 2309; it names a standard fact of the same chapter that the counted argument uses, and every step of the counted proof built on it is retained verbatim in the delivered body.\par\medskip
\ifdefined\hypertarget\hypertarget{sourcepointer-3}{}\fi
\noindent\textit{Source reference.} Original reference 06HJ: sources/stacks-formal-defos/pinned-original-source.tex, source lines 2202--2202. Prerequisite pointer invoked at proof line(s) 2253; it names a standard fact of the same chapter that the counted argument uses, and every step of the counted proof built on it is retained verbatim in the delivered body.\par\medskip
\begin{lemma}
\label{lemma-smooth-morphism-power-series}
Let $R \to S$ be a ring map in $\widehat{\mathcal{C}}_\Lambda$. Then
the induced morphism
$\underline{S}|_{\mathcal{C}_\Lambda} \to \underline{R}|_{\mathcal{C}_\Lambda}$
is smooth if and only if $S$ is a power series ring over $R$.
\end{lemma}

\begin{proof}
Assume $S$ is a power series ring over $R$. Say $S = R[[x_1, \ldots, x_n]]$.
Smoothness of
$\underline{S}|_{\mathcal{C}_\Lambda} \to \underline{R}|_{\mathcal{C}_\Lambda}$
means the following (see Remark \ref{remark-compare-smooth-schlessinger}):
Given a surjective ring map $B \to A$ in
$\mathcal{C}_\Lambda$, a ring map $R \to B$, a ring map $S \to A$ such that
the solid diagram
$$
\xymatrix{
S \ar[r] \ar@{..>}[rd] & A \\
R \ar[u] \ar[r] & B \ar[u]
}
$$
is commutative then a dotted arrow exists making the diagram commute.
(Note the similarity with
Algebra, Definition \ref{algebra-definition-formally-smooth}.)
To construct the dotted arrow choose elements $b_i \in B$ whose images
in $A$ are equal to the images of $x_i$ in $A$. Note that
$b_i \in \mathfrak m_B$ as $x_i$ maps to an element of $\mathfrak m_A$.
Hence there is a unique $R$-algebra map $R[[x_1, \ldots, x_n]] \to B$
which maps $x_i$ to $b_i$ and which can serve as our dotted arrow.

\medskip\noindent
Conversely, assume
$\underline{S}|_{\mathcal{C}_\Lambda} \to \underline{R}|_{\mathcal{C}_\Lambda}$
is smooth. Let $x_1, \ldots, x_n \in S$ be elements whose images
form a basis in the relative cotangent space
$\mathfrak m_S/(\mathfrak m_R S + \mathfrak m_S^2)$ of $S$ over $R$.
Set $T = R[[X_1, \ldots, X_n]]$. Note that both
$$
S/(\mathfrak m_R S + \mathfrak m_S^2) \cong
R/\mathfrak m_R[x_1, \ldots, x_n]/(x_ix_j)
$$
and
$$
T/(\mathfrak m_R T + \mathfrak m_T^2) \cong
R/\mathfrak m_R[X_1, \ldots, X_n]/(X_iX_j).
$$
Let
$S/(\mathfrak m_R S + \mathfrak m_S^2) \to
T/(\mathfrak m_R T + \mathfrak m_T^2)$
be the local $R$-algebra isomorphism given by mapping
the class of $x_i$ to the class of $X_i$. Let
$f_1 : S \to T/(\mathfrak m_R T + \mathfrak m_T^2)$ be the
composition
$S \to S/(\mathfrak m_R S + \mathfrak m_S^2)
\to T/(\mathfrak m_R T + \mathfrak m_T^2)$.
The assumption that
$\underline{S}|_{\mathcal{C}_\Lambda} \to \underline{R}|_{\mathcal{C}_\Lambda}$
is smooth means we can lift $f_1$ to a map
$f_2 : S \to T/\mathfrak{m}_T^2$, then to a map
$f_3 : S \to T/\mathfrak{m}_T^3$, and so on, for all $n \geq 1$. Thus
we get an induced map $f : S \to T = \lim T/\mathfrak m_T^n$
of local $R$-algebras. By our choice of $f_1$, the map $f$ induces an
isomorphism
$\mathfrak m_S/(\mathfrak m_R S + \mathfrak m_S^2) \to
\mathfrak m_T/(\mathfrak m_R T + \mathfrak m_T^2)$
of relative cotangent spaces.
Hence $f$ is surjective by
Lemma \ref{lemma-surjective-cotangent-space}
(where we think of $f$ as a map in $\widehat{\mathcal{C}}_R$).
Choose preimages $y_i \in S$ of $X_i \in T$ under $f$. As $T$ is a
power series ring over $R$ there exists a local
$R$-algebra homomorphism $s : T \to S$ mapping $X_i$ to $y_i$.
By construction $f \circ s = \text{id}$. Then $s$ is injective.
But $s$ induces an isomorphism on relative cotangent spaces since
$f$ does, so it is also surjective by
Lemma \ref{lemma-surjective-cotangent-space}
again. Hence $s$ and $f$ are isomorphisms.
\end{proof}

\end{document}
